% ID: MAT-ALG-INT-007
% Titolo: Numeri interi e proprietà delle potenze 7
% Disciplina: Matematica
% Area: Algebra
% Argomento: Numeri interi
% Classe: Prima secondaria di secondo grado
% Difficolta: 3
% Tipo: problema_numerico
% Tag: matematica, algebra, numeri_interi
% Fonte: verifica_interi_anonimizzata
% Provenienza: materiale fornito dal docente; testo anonimizzato; soluzione redatta per Exergo

\begin{esercizio}
Tre amici decidono di configurare il loro computer personale in modo che esegua automaticamente il backup dei dati. Anna sceglie di salvare i dati ogni settimana, Barbara ogni quattro giorni, Andrea ogni due settimane. Se il primo salvataggio automatico avviene per tutti venerdì 9 agosto, in quale giorno si verificherà di nuovo l'evento? In quale giorno avverrà il primo salvataggio per Andrea e Anna ma non per Barbara?
\end{esercizio}

\begin{soluzione}
I periodi sono 7, 4 e 14 giorni. Il minimo comune multiplo è 28:
la prima coincidenza successiva avviene venerdì 6 settembre.
Anna e Andrea coincidono ogni 14 giorni; la prima volta, venerdì 23 agosto,
non coincide con il periodo di 4 giorni di Barbara.
\end{soluzione}
